\scanpage{391}
\mspage{198}
se remonte de façon unique en
$\unsure{\widehat{\widehat{\mathcal G}}^{G}}$.
Il faudrait réexaminer ceci de façon plus \unsure{soigneuse} par la suite.
Mais il me semble qu'on ne trouve pas ici des nouvelles propriétés des
$\Gamma\subset\unsure{\mathcal M_a|\mathcal N_a}$ ou de
$\unsure{M_a|N_a}$ lui $\overline{\mathrm m}$, i.e. de
$\unsure{\Pi_a}$ comme ss-groupe profini de
$\Gamma'=\operatorname{Norm}_{\unsure{\widehat{\widehat{\mathcal T}}}}
(\unsure{\widehat{\mathcal T}})/\unsure{\widehat{\mathcal T}}$

\uline{31}\quad Digression sur les \unsure{relèvements}
d'\inserted{une} \unsure{action} extérieure
d\deleted{\unsure{e}}\inserted{'un} groupe fini $G$ sur un
groupe profini à \unsure{bouts} $\pi$.
\figure{[Trait vertical à gauche des deux dernières lignes du titre.]}

On suppose l'action extérieure \uline{fidèle},
i.e. $G\subset\unsure{\widehat{\widehat{\mathcal T}}}(\pi)
=\unsure{\widehat{\widehat{\mathcal T}}}$,
\inserted{et que $\unsure{G=G^+}$.}
On suppose de plus qu'il existe une \unsure{discrétification}
invariante $\pi_0$, i.e. telle que
$G\subset\unsure{\mathcal T}(\pi_0)$.

Si on suppose $G$ cyclique, dans ce \unsure{cas}
encore il est \unsure{prouvé} que la situation est
\uline{visible} topologiquement (\unsure{donc} :
\uline{\unsure{algébr.}} de façon analytique \unsure{complexe} $\ldots$).

\scanpage{392}
Dans le cas discret, la signification des classes de
\inserted{$\pi_0$-}conjugaison de \unsure{relèvements}
$\unsure{\widehat G|G}\longrightarrow\unsure{\mathcal G}(\pi_0)$ est bien comprise ;
\deleted{Mais}\inserted{de $\overline{\mathrm m}$ que pour les relèvements partiels}
On trouve une \deleted{\unsure{cat.} $U^{!}$} réalisation
\inserted{\unsure{canonique}} du groupoïde fondamental
\margin{si $U^{!}\ne\unsure{\varnothing|\phi}$}
associé au
groupe ext. $\pi_0$, par un groupoïde \inserted{fini} $[U^{!}]$
sur lequel $\unsure{\widehat G|G}$ opère au sens strict, dont les pts
correspondent aux classes de $\pi_0$-conjugaison de sections
partielles $\ne1$ évidemment, le groupe
\deleted{\illegible{1}}$\unsure{\mathcal T}(\pi_0)^G$
opère de façon \unsure{rigide} canonique sur ce groupoïde.
Si $P\in U^{!}$ correspond à un \unsure{relèvement}
$\unsure{\widehat G|G}\longrightarrow\unsure{\mathcal G}(\pi_0)$, alors
un élément $\deleted{\unsure{g}}\inserted{\dot u}\in
\unsure{\mathcal T}(\pi_0)^G$ est dans l'image de
$\unsure{\mathcal G}(\pi_0)^G$, i.e. peut se remonter en
$u\in\unsure{\mathcal G}(\pi_0)$ commutant à $\unsure{\widehat G|G}$, sss
$\dot u(P)=P$ [i.e. sss
\inserted{(si $\dot u$ est remonté de façon \unsure{quelc.} en $u$)}
$\unsure{u}(\unsure{\widehat G|G})$ est $\pi_0$-conjugué à $\unsure{\widehat G|G}$
i.e. sss $\exists g\in\pi_0$ tel que
$\unsure{u}(\unsure{\widehat G|G})=\operatorname{int}(g)(\unsure{\widehat G|G})$
i.e. \deleted{\unsure{qui}}\inserted{alors}
$\unsure{u'}=\operatorname{int}(g^{-1})u$ fixe $\unsure{\widehat G|G}$
\inserted{(par l'hyp. $\dot u\in
\unsure{\widehat{\mathcal T}^{G}|\mathcal T^{G}}$)}.
\deleted{\illegible{1}} Comme on \unsure{devait} \unsure{s'y}
\unsure{attendre!} toute notre belle assertion, est une tautologie]

\scanpage{393}
\mspage{199}
et donc le \unsure{relèvement} est unique. Comme
$U^{!}_{\unsure{*|\varphi}}$ est fini, donc le stabilisateur
$\unsure{\mathcal T}(\pi_0)^G_P$ de $P$ dans
$\unsure{\mathcal T}(\pi_0)^G$ est d'indice fini,
et c'est donc un ss-groupe d'indice fini qui se remonte
globalement et canoniquement.

Que peut-on dire dans le contexte profini ? Bien sur,
on a une application canonique
\diagram{\begin{tikzcd}[ampersand replacement=\&,column sep=large]
\begin{gathered}
\text{\deleted{\illegible{1}} classes de $\pi_0$-conj.}\\
\text{de relèvements de}\\
\unsure{G|\widehat G}\text{ en }\unsure{G|\widehat G}
\longrightarrow\unsure{\mathcal G}(\pi_0)\\
\text{i.e. de scindages de}\\
\text{l'ext. }1\longrightarrow\pi_0\longrightarrow
\unsure{E|\mathcal E}\longrightarrow G\longrightarrow1
\end{gathered}
\arrow[r] \&
\begin{gathered}
\text{classes de $\unsure{\pi_0|\pi}$-}\\
\text{conj. de \unsure{rel.} de}\\
\unsure{G|\widehat G}\text{ en }\unsure{G|\widehat G}
\longrightarrow\unsure{\widehat{\widehat{\mathcal G}}}(\pi_0)\\
\text{i.e. de scindages}\\
\text{de l'extension}\\
1\longrightarrow\underset{\substack{\Vert\\\pi}}{\widehat{\pi}_0}
\longrightarrow\unsure{\widehat E|\widehat{\mathcal E}}
\longrightarrow G\longrightarrow1
\end{gathered}
\end{tikzcd}}

Il faudrait examiner d'abord

a) La question de la bijectivité de cette appl.

b) La question si $\pi^G=1$ (\unsure{\uline{rigidité}|rigidité}),
pour $G\ne(1)$. -- (pour un relèvement donné,
dans $\unsure{E|\mathcal E}$ pour simplifier)

J'ai envie de conjecturer \unsure{sans vergogne}
qu'il en est ainsi --- ce qui impliquerait
p.ex. que pour tt tel \unsure{relèvement} $\unsure{\widehat G|G}$, il y a

\scanpage{394}

un sous-groupe \deleted{\illegible{1}} \inserted{$(\widehat{\widehat{\unsure{\mathcal T|\tau}}}^{G})_{P}$} du groupe (peut-être vraiment \unsure{immense} à priori !) $\widehat{\widehat{\unsure{\mathcal T|\tau}}}^{G}$, qui se remonte \unsure{canoniquement} de façon à \unsure{commuter} à \unsure{$\widetilde G$|$G$}....

Il faudrait manifestement faire, en $\overline{\mathrm m}$ temps qu'une théorie des op. ext. de groupes finis sur des groupes discrets à \unsure{bouts|lacets} (qui est pour le moment entièrement conjecturale) une théorie analogue dans le cas profini --- j'aurai sans doute à y revenir par la suite.

\scanpage{395}
\mspage{200}

\uline{32} Retour sur les aspects arithmétiques des \unsure{bouchages} de trous : relations entre $\Gamma_{g,\nu}$ et $\Gamma_{g,\nu-1}$.

Bien sûr, on \deleted{\unsure{voit que}} a que
\[
\mathcal N_{g,\nu}\longrightarrow\unsure{\mathfrak S|G}_{\nu}
\]
est surjectif \inserted{(puisque $\unsure{\mathcal T|\tau}_{g,\nu}\longrightarrow\unsure{\mathfrak S|G}_{\nu}$ l'est)} \unsure{donc} on a
\[
(1)\qquad 1\unsure{\longrightarrow|-}\mathcal N_{g,\nu}^{!}\unsure{\longrightarrow|-}\mathcal N_{g,\nu}\unsure{\longrightarrow|-}\unsure{\mathfrak S|G}_{\nu}\unsure{\longrightarrow|-}1
\]
et le diagramme cartésien de \unsure{ssgroupes} de $\widehat{\widehat{\unsure{\mathcal T|\tau}}}_{g,\nu}$
\[
(2)\qquad
\diagram{\begin{tikzcd}[ampersand replacement=\&]
\widehat{\unsure{\mathcal T|\tau}}_{g,\nu}^{+}\arrow[d] \& \mathcal N_{g,\nu}\arrow[l]\arrow[d]\\
\widehat{\unsure{\mathcal T|\tau}}_{g,\nu}^{!+} \& \mathcal N_{g,\nu}^{!}\arrow[l]
\end{tikzcd}}
\]
\margin{\deleted{\illegible{1}} $\mathcal N_{g,\nu}=\mathcal N_{g,\nu}^{!}\cdot\widehat{\unsure{\mathcal T|\tau}}_{g,\nu}^{+}$ i.e. $\mathcal N_{g,\nu}$ \unsure{connu} par $\mathcal N_{g,\nu}^{!}$}
\unsure{donne|donnera} un isomorphisme
\[
(3)\qquad \mathcal N_{g,\nu}^{!}/\widehat{\unsure{\mathcal T|\tau}}_{g,\nu}^{!+}\xrightarrow{\sim}\mathcal N_{g,\nu}/\widehat{\unsure{\mathcal T|\tau}}_{g,\nu}^{+}=\mathbb{\Gamma}_{g,\nu}
\]
Dans le cas d'un $(\unsure{\Pi|\pi},a)$, \inserted{($\unsure{\Pi|\pi}$ groupe ext. à \unsure{bouts|lacets},} $a$ une arithmétisation), on aura de $\overline{\mathrm m}$
\[
(4)\qquad\left\{\begin{aligned}
\mathcal N_a^{!}/\widehat{\unsure{\mathcal T|\tau}}_a^{+}&\simeq\mathcal N_a/\widehat{\unsure{\mathcal T|\tau}}_a^{+}=\mathbb{\Gamma}_a.\\
\mathcal N_a&=\mathcal N_a^{!}\widehat{\unsure{\mathcal T|\tau}}_a^{+}
\end{aligned}\right.
\]

Considérons maintenant \deleted{une \unsure{partie}\inserted{\unsure{$i\in I=I(\Pi)$}} , \illegible{3} $\unsure{(g,\nu)}$, \unsure{et un} $\unsure{\Pi_i'|\pi_i'}$ \illegible{2} $\unsure{I=I(\Pi)}$, et soit \illegible{1}} pour $\forall i\in I$ le stabilisateur $\mathcal N_i$ de $i$ dans $\mathcal N$, et le groupe \deleted{\unsure{extérieur}} \inserted{à \unsure{bouts|lacets} (gr. ext.!)} $\unsure{\Pi_i'|\pi_i'}$ de type $g,\nu-1$, déduit de $\unsure{\Pi|\pi}$ par ``bouchage du trou $i$'',
\margin{On écrit $\mathcal N$ etc. au lieu de $\mathcal N_a$. On \unsure{suppose} \illegible{2} $(g,\nu-1)$, i.e. \unsure{$2g-2+\nu$} \illegible{1} \unsure{$\geq1$}, i.e. \unsure{$2g+\nu\geq4$}}

\scanpage{396}

on a alors, pour une \unsure{discrétification} choisie, \deleted{\illegible{1}} \inserted{$\unsure{\alpha|a}$} compatible avec l'arithmétisation, un \unsure{hom.|hom} injectif
\[
(5)\qquad \unsure{\Pi_i'|\pi_i'}\xrightarrow{\varphi_i}\widehat{\unsure{\mathcal T|\tau}}_{\unsure{i|a}}^{+!}\subset\mathcal N_{\unsure{i|a}}^{!}\subset\widehat{\widehat{\unsure{\mathcal T|\tau}}}_i
\]

\deleted{On \unsure{ne peut} \illegible{2}}
\deleted{On aura}
\deleted{\unsure{qui remonte le composé}}
\deleted{\textcircled{\unsure{e|c}} L'image de $\unsure{\Pi'|\pi'}$}
tel que le composé
\[
\diagram{\begin{tikzcd}[ampersand replacement=\&]
\unsure{\Pi_i'|\pi_i'}\arrow[r] \& \widehat{\widehat{\unsure{\mathcal T|\tau}}}_i\arrow[r] \& \widehat{\widehat{\unsure{G|\mathcal G}}}(\unsure{\Pi_i'|\pi_i'})\\
\&\& \unsure{\Pi_i'|\pi_i'}\arrow[u,hook]
\end{tikzcd}}
\]
soit l'inclusion canonique. Ceci \unsure{posé}, on veut\margin{plutôt : \uline{il est vrai que !}} que l'image de (5) (qui n'est peut-être pas invariante dans $\widehat{\widehat{\unsure{\mathcal T|\tau}}}_i$) soit \uline{invariante dans \deleted{$\unsure{\mathcal T|\tau}$} $\mathcal N_i$}, ou ce qui revient au $\overline{\mathrm m}$, dans $\mathcal N^{!}$ (car $\mathcal N_i=\mathcal N^{!}\cdot\widehat{\unsure{\mathcal T|\tau}}_i^{+}$ or $\widehat{\unsure{\mathcal T|\tau}}_i^{+}$ \unsure{invarie} cette image). \deleted{\illegible{1}} Revenant à la situation universelle, cela \unsure{\uline{signifie}|signifie}
\margin{Cela signifie aussi que $\varphi_i$ \unsure{ne dépend pas} du choix de la \unsure{discrétification} \unsure{de donnée} $\alpha$...}

\textcircled{\unsure{e|c}} $\mathcal N_{g,\nu}^{!}\subset\widehat{\widehat{\unsure{\mathcal T|\tau}}}_{g,\nu}$ est contenu dans le normalisateur des $\unsure{\Pi_i'|\pi_i'}\hookrightarrow\widehat{\unsure{\mathcal T|\tau}}_{g,\nu}^{!+}$ $(0\leq i\leq\nu-1)$ \inserted{qui sont donc invariants dans $\mathcal N_{g,\nu}^{!}$} \deleted{\unsure{de} $(g,\nu-1)$ \illegible{2}}.

De plus $\widehat{\widehat{\unsure{\mathcal T|\tau}}}_{g,\nu}^{!}\xrightarrow{\psi_i}\widehat{\widehat{\unsure{G|\mathcal G}}}_{g,\nu-1}^{!}$ induit un \deleted{\illegible{1}} \uuline{\unsure{iso|isom.}}
\[
(6)\qquad\mathcal N_{g,\nu}^{!}\xrightarrow{\sim}\unsure{\mathcal M|\mathcal N}_{g,\nu-1}^{!}
\]
\margin{\uline{NB} Il suffit \unsure{un seul} \illegible{1} \unsure{dans les} \illegible{2} pour \illegible{2} \deleted{\illegible{1}}}
\margin{\uline{NB} \illegible{4}}

\scanpage{397}
\mspage{201}

Cette assertion se décompose en deux : \unsure{donc|tout} d'abord que $\psi_i$ applique bien $\mathcal N_{g,\nu}^{!}$ dans $\mathcal M_{g,\nu-1}^{!}$ -- et ceci résulte de la définition \unsure{arithmético-géom} de ces groupes -- cela implique \unsure{d'autre} part, puisque $\psi_i$ induit aussi un \inserted{\unsure{iso}} \deleted{\illegible{1}}
\[
\unsure{\widehat{\mathcal T}_{g,\nu}^{+}}
\xrightarrow{\sim}
\unsure{\widehat{\mathcal G}_{g,\nu-1}^{+}},
\]
qu'il induit un hom
\[
(7)\qquad
\begin{array}{ccc}
\unsure{\mathbb{\Gamma}}_{g,\nu}
&\xrightarrow{\substack{\inserted{\lambda_i}\\\deleted{\unsure{\psi_i}}}}&
\unsure{\mathbb{\Gamma}}_{g,\nu-1}\\
\Vert&&\Vert\\
\mathcal N_{g,\nu}/\unsure{\widehat{\mathcal T}_{g,\nu}^{+}}
&&\mathcal M_{g,\nu-1}/\unsure{\widehat{\mathcal G}}
\end{array}
\]
\deleted{et par voie \unsure{arithmétique}}
et l'inj. (resp. bijectivité) de (6) équivaut à celle de (7). D'ailleurs (7) s'insère, par construction, dans un diagramme commutatif

(8)\quad\diagram{\begin{tikzcd}[ampersand replacement=\&,column sep=large,row sep=large]
\& \unsure{\mathbb{\Gamma}}_{\mathbb Q}
  \arrow[dl,"\Theta_{g,\nu}"']
  \arrow[dr,"\Theta_{g,\nu-1}"] \& \\
\unsure{\mathbb{\Gamma}}_{g,\nu} \arrow[rr]
\&\&\unsure{\mathbb{\Gamma}}_{g,\nu-1}
\end{tikzcd}}

avec $\Theta_{g,\nu}$ et $\unsure{\Theta_{g,\nu-1}|\Theta_{g,\nu+1}}$ surjectifs, donc il est bel et bien évident que (6) est \uline{surjectif}. La bijectivité \unsure{\uline{signifie}|signifie} que $\Theta_{g,\nu}$ et $\Theta_{g,\nu-1}$ ont $\overline{\mathrm m}$ noyau \unsure{(alors} \unsure{qu'à priori} il se \unsure{pourrait} que $\Theta_{g,\nu}$ ait un noyau plus petit i.e. corresponde à une représentation de $\unsure{\mathbb{\Gamma}}_{\mathbb Q}$ ``moins infidèle'' que $\Theta_{g,\nu-1}$. D'ailleurs,

\scanpage{398}

il est évident \deleted{que} ($\unsure{\overline{\mathrm m}}$ à priori, sans utiliser la définition \inserted{\unsure{arithm}} explicite des $\mathcal N_{g,\nu}$) que l'hom $\deleted{\illegible{1}}\inserted{\unsure{\lambda_i}}$ \inserted{\unsure{(7)}} ne dépend pas du choix de $0\leq i\leq\nu-1$ -- p.ex. puisque l'on \unsure{passe} de l'un à l'autre en appliquant des opérations de $\unsure{\widehat{\mathcal T}^{+}}$ (puisque $\unsure{\widehat{\mathcal T}^{+}}$ opère \unsure{transitivement} sur $I$) et que $\unsure{\widehat{\mathcal T}^{+}}$ opère \unsure{trivialt|trivialement} dans $\unsure{\mathbb{\Gamma}_{g,\nu}}$ ...

On peut dire que (6) décrit $\mathcal M_{g,\nu-1}^{!}$ (donc $\mathcal N_{g,\nu-1}^{!}$) en termes de $\mathcal N_{g,\nu}$ -- mais \unsure{l'inverse|l'inconnu} est moins \unsure{clair}, \unsure{faute} de savoir si
\[
\unsure{\widehat{\widehat{\mathcal T}}_{g,\nu}^{!}}
\xrightarrow{\psi_i}
\unsure{\widehat{\widehat{\mathcal G}}_{g,\nu-1}}
\]
est injectif ; \unsure{si on} le savait, on pourrait décrire $\mathcal N_{g,\nu}^{!}$ comme l'image inverse de $\mathcal M_{g,\nu-1}^{!}$ ... \deleted{\unsure{Alors}} Il ne serait pas impossible d'ailleurs que (f) soit faux, i.e. que \inserted{les} $\Theta_{g,\nu}$ soient infidèles, mais de moins en moins quand on fait augmenter $\nu$ -- en passant à la limite \inserted{\unsure{projective}} $\Theta_{g,\infty}$, \unsure{du moins}, aurait-on

\scanpage{399}
\mspage{202}

(peut-être !) une \unsure{représentation} fidèle de $\unsure{\mathbb{\Gamma}}_{\mathbb Q}$. Mais jusqu'à indication contraire, je \unsure{préfère} travailler \unsure{hypothétiquement} avec (f), ce qui s'exprime par les suites \unsure{exactes|ex.tes} \unsure{fondamentales}
\[
(9)\qquad
1\longrightarrow\widehat{\Pi}_{g,\nu-1}
\xrightarrow{\varphi_i}
\unsure{\mathcal N_{g,\nu}^{!}|\mathcal M_{g,\nu}^{!}}
\unsure{\overset{\psi_i'}{\longleftrightarrow}|\xrightarrow{\psi_i'}}
\unsure{\mathcal N_{g,\nu-1}^{!}|\mathcal M_{g,\nu-1}^{!}}
\longrightarrow1
\]
\margin{\unsure{\&} on pourrait l'inclure dans une suite \unsure{ex.te} un peu plus grande, avec $(\mathcal N_{g,\nu})_i$ et $\mathcal N_{g,\nu-1}$ ...}
qui étend la suite \unsure{ex.te|exacte}
\[
(10)\qquad
1\longrightarrow\widehat{\Pi}_{g,\nu-1}
\xrightarrow{\varphi_i}
\unsure{\widehat{\mathcal T}_{g,\nu}^{!}}
\xrightarrow{\psi_i'}
\unsure{\widehat{\mathcal T}_{g,\nu-1}^{!}}
\longrightarrow1
\]
(et \unsure{tient|tiens} lieu de la suite \unsure{exacte|ex.te} peut-être \unsure{défaillante}
\[
??\qquad
1\longrightarrow\unsure{\widehat{\Pi}_{g,\nu-1}|\widehat{\pi}_{g,\nu-1}}
\longrightarrow\unsure{\widehat{\widehat{\mathcal T}}_{g,\nu}^{!}}
\longrightarrow\unsure{\widehat{\widehat{\mathcal T}}_{g,\nu-1}^{!}}
\longrightarrow1
\qquad??
\]
)

Quand $(\unsure{\Pi|\pi},a)$ est un groupe \unsure{extérieur} à \unsure{bouts|bords} muni d'une \unsure{arithmétisation} -- ce qu'on pourrait appeler une ``courbe algébrique virtuelle'' -- alors ce qui précède permet de définir, sur \inserted{chacun} des $\unsure{\Pi_i'|\pi_i'}$ de type $(g,\nu-1)$ associés aux $i\in\unsure{I|\Sigma}(\unsure{\Pi|\pi})$, \unsure{une} \unsure{arithmétisation} \unsure{canoniquement} associée à $a$, soit $\unsure{a'}$,

\scanpage{400}

et on \unsure{trouve} \unsure{donc} une suite \unsure{exacte|exact}
\[
1 \longrightarrow \pi'_i \longrightarrow (\mathcal{N}_a)_i
\xrightarrow{\psi_i} \mathcal{N}_{a'_i} \longrightarrow 1
\tag{11}
\]
telle que l'on ait
\[
(\widehat{\unsure{\mathcal{T}|\tau}}^{+}_a)_i
=\psi_i^{-1}(\widehat{\unsure{\mathcal{T}|\tau}}^{+}_{a'_i}),
\tag{12}
\]
\[
\downarrow
\]
\unsure{induisant}

\diagram{[(11') : une ligne horizontale
$1\longrightarrow\pi'_i\longrightarrow\mathcal{N}_{\unsure{a}}^{!}
\xrightarrow{\psi_i^{!}}\mathcal{N}_{a'_i}^{!}\longrightarrow1$.
Sous $\mathcal{N}_{\unsure{a}}^{!}$, un signe vertical
$\unsure{\cup|\downarrow}$, puis
$\widehat{\unsure{\mathcal{T}|\tau}}_a^{+!}$ ; sous ce dernier,
un signe $=$ écrit verticalement, puis
$\unsure{\psi_i^{\prime-1}|\psi_i^{-1}}(\widehat{\unsure{\mathcal{T}|\tau}}_{\unsure{a_i}}^{+!}$.
La parenthèse fermante n'est pas lisible.
Sous $\mathcal{N}_{a'_i}^{!}$, un signe vertical
$\unsure{\cup|\downarrow}$, puis
$\widehat{\unsure{\mathcal{T}|\tau}}_{a'_i}^{+!}$.]}

et \unsure{induisant} par \unsure{passage} aux quotients
\[
\unsure{\mathbb{\Gamma}}_a\xrightarrow{\sim}
\unsure{\mathbb{\Gamma}}_{\unsure{a_i}}
\tag{13}
\]
De plus, \uline{\unsure{l'application}} canonique
\[
\unsure{\mathrm{Discr}}^{+}(\pi)\longrightarrow
\unsure{\mathrm{Discr}}^{+}(\pi'_i)
\]
\unsure{définit} par \unsure{passage} aux quotients
\[
\mathbb{P}_a\xrightarrow{\sim}\mathbb{P}_{\unsure{a_i}}
\tag{14}
\]
compatible avec les actions de $\unsure{\mathbb{\Gamma}}_a$,
$\unsure{\mathbb{\Gamma}}_{a'_i}$ et (13).

Je \unsure{me} fais \inserted{ici} encore \unsure{assertion|assertions} sur
une \unsure{soit-disante|soi-disante} bijectivité entre
ens. des arithmétisations de $\pi$, et \unsure{ens} des
